% Appendix C — DE-MLP loss, ensemble aggregation, and training configuration.
% Moved out of Methods Sect. 3.3 (predictors + probabilistic model) to keep the
% main text to predictor groups and essential architecture.
%
% Source of truth:
%   src/models/deep_ensemble.py   (GaussianMLP, beta_nll, ensemble_predict, main)
%   src/models/train_common.py    (TrainConfig: AdamW/OneCycle/grad-clip/early-stop)
%   src/models/pipeline.py        (feature sets, target columns)
%
% Usage: % Appendix C — DE-MLP loss, ensemble aggregation, and training configuration.
% Moved out of Methods Sect. 3.3 (predictors + probabilistic model) to keep the
% main text to predictor groups and essential architecture.
%
% Source of truth:
%   src/models/deep_ensemble.py   (GaussianMLP, beta_nll, ensemble_predict, main)
%   src/models/train_common.py    (TrainConfig: AdamW/OneCycle/grad-clip/early-stop)
%   src/models/pipeline.py        (feature sets, target columns)
%
% Usage: % Appendix C — DE-MLP loss, ensemble aggregation, and training configuration.
% Moved out of Methods Sect. 3.3 (predictors + probabilistic model) to keep the
% main text to predictor groups and essential architecture.
%
% Source of truth:
%   src/models/deep_ensemble.py   (GaussianMLP, beta_nll, ensemble_predict, main)
%   src/models/train_common.py    (TrainConfig: AdamW/OneCycle/grad-clip/early-stop)
%   src/models/pipeline.py        (feature sets, target columns)
%
% Usage: % Appendix C — DE-MLP loss, ensemble aggregation, and training configuration.
% Moved out of Methods Sect. 3.3 (predictors + probabilistic model) to keep the
% main text to predictor groups and essential architecture.
%
% Source of truth:
%   src/models/deep_ensemble.py   (GaussianMLP, beta_nll, ensemble_predict, main)
%   src/models/train_common.py    (TrainConfig: AdamW/OneCycle/grad-clip/early-stop)
%   src/models/pipeline.py        (feature sets, target columns)
%
% Usage: \input{MANUSCRIPT_APPENDIX_MODEL.tex}
% Requires in the main preamble:
%   \usepackage{amsmath}
%   \usepackage{booktabs}
%   \usepackage{natbib}   % \citet  (or map to your journal's citation macro)
%
% NOTE: batch size / max-epoch values below are the production (Linux/HPC) run
% settings; the macOS quick-iteration defaults (2048 / 100) are not used for the
% frozen models.

\section{Deep-ensemble loss, aggregation, and training configuration}
\label{app:model}

\subsection{Member architecture and probabilistic head}

Each of the $M = 5$ ensemble members is a Gaussian-head multilayer perceptron
with two hidden layers (64 and 32 units). Every hidden block applies, in order,
a linear map, layer normalization, a ReLU activation, and 10\% dropout; dropout
is active during training only, so there is no Monte-Carlo dropout at inference
and the ensemble spread alone provides the epistemic uncertainty. For footprint
$i$, member $m$ outputs a predicted mean $\mu_{m,i}$ and a raw second output that
is clamped to $[-10, 10]$ and interpreted as the log-variance
$s_{m,i} = \log\sigma_{m,i}^{2}$, giving a per-footprint predictive variance
$\sigma_{m,i}^{2} = \exp(s_{m,i})$.

\subsection{Training loss}

Writing the target as $y_i \equiv \Delta \mathrm{X_{CO2,\,i}^{B11}}$, the standard
Gaussian negative log-likelihood of one footprint is, up to an additive constant,
%
\begin{equation}
  \ell^{\mathrm{NLL}}_{m,i}
  = \frac{1}{2}\left[\,s_{m,i}
  + \frac{(y_i-\mu_{m,i})^{2}}{\sigma_{m,i}^{2}}\,\right].
  \label{eq:gauss-nll}
\end{equation}
%
Because this loss down-weights high-variance footprints, the mean is poorly
fitted exactly in the near-cloud regime where $\sigma_{m,i}^{2}$ is large. We
therefore train the members with the $\beta$-NLL loss of
\citet{seitzer2022pitfalls}, which multiplies each footprint's NLL by a
stop-gradient function of its own predicted variance,
%
\begin{equation}
  \ell^{\beta\text{-}\mathrm{NLL}}_{m,i}
  = \bigl\lfloor\sigma_{m,i}^{2}\bigr\rfloor_{\mathrm{sg}}^{\,\beta}\;
  \ell^{\mathrm{NLL}}_{m,i},
  \qquad \beta = 0.5,
  \label{eq:beta-nll}
\end{equation}
%
where $\lfloor\cdot\rfloor_{\mathrm{sg}}$ denotes a stop-gradient (the weight is
treated as a constant during backpropagation). Here $\beta = 0$ recovers the
plain NLL of Eq.~\eqref{eq:gauss-nll} and $\beta = 1$ recovers a
mean-squared-error weighting of the mean; the production choice $\beta = 0.5$
restores the mean-fitting gradient in the high-variance near-cloud footprints
while keeping a calibrated variance. The member loss is the average of
Eq.~\eqref{eq:beta-nll} over the training footprints.

\subsection{Ensemble aggregation}

The five members differ only in their random seed (member $m$ uses seed
$s_0 + m$), so their predictions decorrelate and the ensemble spread provides
the epistemic uncertainty. The per-member means and variances are combined as a
Gaussian mixture,
%
\begin{equation}
  \mu^{*}_i = \frac{1}{M}\sum_{m=1}^{M}\mu_{m,i},
  \qquad
  \sigma^{*2}_i
  = \underbrace{\frac{1}{M}\sum_{m=1}^{M}\sigma_{m,i}^{2}}_{\text{aleatoric}}
  + \underbrace{\frac{1}{M}\sum_{m=1}^{M}\mu_{m,i}^{2} - \mu^{*2}_i}_{\text{epistemic}},
  \label{eq:mixture}
\end{equation}
%
which is the law of total variance: the first term is the mean aleatoric
variance predicted by the heads and the second is the between-member (epistemic)
variance. The ensemble mean $\mu^{*}_i$ is the point correction applied to each
footprint.

\subsection{Optimization and leakage control}

Each member is optimized with AdamW (learning rate $10^{-3}$, weight decay
$10^{-4}$) under a one-cycle learning-rate schedule with gradient-norm clipping
at $1.0$, a batch size of $4096$, and up to $500$ epochs. Training is stopped
early on the held-out calibration block (patience of $50$ epochs), and the
checkpoint with the lowest calibration loss is kept; the calibration block is
carved from the training dates and never overlaps the held-out fold or the
validation references. The calibration block is also used to recalibrate the
Mondrian conformal intervals, which are binned by predicted-anomaly deciles.
Feature scaling and any profile-EOF transforms are fitted on the proper-training
partition only. Table~\ref{tab:hyperparams} collects the full configuration.

\begin{table}[t]
\centering
\caption{DE-MLP architecture and training configuration (production settings).}
\label{tab:hyperparams}
\begin{tabular}{@{}ll@{}}
\toprule
Component & Setting \\
\midrule
Ensemble members $M$        & 5 (seeds $s_0 + m$) \\
Hidden layers               & 64, 32 \\
Hidden block                & Linear $\to$ LayerNorm $\to$ ReLU $\to$ Dropout(0.1) \\
Output head                 & Linear $\to$ 2 (mean, log-variance) \\
Log-variance clamp          & $[-10, 10]$ \\
Loss                        & $\beta$-NLL, $\beta = 0.5$ \\
Optimizer                   & AdamW \\
Learning rate (max)         & $10^{-3}$ \\
Weight decay                & $10^{-4}$ \\
LR schedule                 & one-cycle (\texttt{pct\_start} 0.05, div 25, final div 1000) \\
Gradient clipping           & 1.0 (max norm) \\
Batch size                  & 4096 \\
Maximum epochs              & 500 \\
Early stopping              & patience 50, on the calibration block \\
Calibration fraction        & 0.15 of the training dates \\
Conformal intervals         & Mondrian, 10 predicted-anomaly decile bins \\
\bottomrule
\end{tabular}
\end{table}

% Requires in the main preamble:
%   \usepackage{amsmath}
%   \usepackage{booktabs}
%   \usepackage{natbib}   % \citet  (or map to your journal's citation macro)
%
% NOTE: batch size / max-epoch values below are the production (Linux/HPC) run
% settings; the macOS quick-iteration defaults (2048 / 100) are not used for the
% frozen models.

\section{Deep-ensemble loss, aggregation, and training configuration}
\label{app:model}

\subsection{Member architecture and probabilistic head}

Each of the $M = 5$ ensemble members is a Gaussian-head multilayer perceptron
with two hidden layers (64 and 32 units). Every hidden block applies, in order,
a linear map, layer normalization, a ReLU activation, and 10\% dropout; dropout
is active during training only, so there is no Monte-Carlo dropout at inference
and the ensemble spread alone provides the epistemic uncertainty. For footprint
$i$, member $m$ outputs a predicted mean $\mu_{m,i}$ and a raw second output that
is clamped to $[-10, 10]$ and interpreted as the log-variance
$s_{m,i} = \log\sigma_{m,i}^{2}$, giving a per-footprint predictive variance
$\sigma_{m,i}^{2} = \exp(s_{m,i})$.

\subsection{Training loss}

Writing the target as $y_i \equiv \Delta \mathrm{X_{CO2,\,i}^{B11}}$, the standard
Gaussian negative log-likelihood of one footprint is, up to an additive constant,
%
\begin{equation}
  \ell^{\mathrm{NLL}}_{m,i}
  = \frac{1}{2}\left[\,s_{m,i}
  + \frac{(y_i-\mu_{m,i})^{2}}{\sigma_{m,i}^{2}}\,\right].
  \label{eq:gauss-nll}
\end{equation}
%
Because this loss down-weights high-variance footprints, the mean is poorly
fitted exactly in the near-cloud regime where $\sigma_{m,i}^{2}$ is large. We
therefore train the members with the $\beta$-NLL loss of
\citet{seitzer2022pitfalls}, which multiplies each footprint's NLL by a
stop-gradient function of its own predicted variance,
%
\begin{equation}
  \ell^{\beta\text{-}\mathrm{NLL}}_{m,i}
  = \bigl\lfloor\sigma_{m,i}^{2}\bigr\rfloor_{\mathrm{sg}}^{\,\beta}\;
  \ell^{\mathrm{NLL}}_{m,i},
  \qquad \beta = 0.5,
  \label{eq:beta-nll}
\end{equation}
%
where $\lfloor\cdot\rfloor_{\mathrm{sg}}$ denotes a stop-gradient (the weight is
treated as a constant during backpropagation). Here $\beta = 0$ recovers the
plain NLL of Eq.~\eqref{eq:gauss-nll} and $\beta = 1$ recovers a
mean-squared-error weighting of the mean; the production choice $\beta = 0.5$
restores the mean-fitting gradient in the high-variance near-cloud footprints
while keeping a calibrated variance. The member loss is the average of
Eq.~\eqref{eq:beta-nll} over the training footprints.

\subsection{Ensemble aggregation}

The five members differ only in their random seed (member $m$ uses seed
$s_0 + m$), so their predictions decorrelate and the ensemble spread provides
the epistemic uncertainty. The per-member means and variances are combined as a
Gaussian mixture,
%
\begin{equation}
  \mu^{*}_i = \frac{1}{M}\sum_{m=1}^{M}\mu_{m,i},
  \qquad
  \sigma^{*2}_i
  = \underbrace{\frac{1}{M}\sum_{m=1}^{M}\sigma_{m,i}^{2}}_{\text{aleatoric}}
  + \underbrace{\frac{1}{M}\sum_{m=1}^{M}\mu_{m,i}^{2} - \mu^{*2}_i}_{\text{epistemic}},
  \label{eq:mixture}
\end{equation}
%
which is the law of total variance: the first term is the mean aleatoric
variance predicted by the heads and the second is the between-member (epistemic)
variance. The ensemble mean $\mu^{*}_i$ is the point correction applied to each
footprint.

\subsection{Optimization and leakage control}

Each member is optimized with AdamW (learning rate $10^{-3}$, weight decay
$10^{-4}$) under a one-cycle learning-rate schedule with gradient-norm clipping
at $1.0$, a batch size of $4096$, and up to $500$ epochs. Training is stopped
early on the held-out calibration block (patience of $50$ epochs), and the
checkpoint with the lowest calibration loss is kept; the calibration block is
carved from the training dates and never overlaps the held-out fold or the
validation references. The calibration block is also used to recalibrate the
Mondrian conformal intervals, which are binned by predicted-anomaly deciles.
Feature scaling and any profile-EOF transforms are fitted on the proper-training
partition only. Table~\ref{tab:hyperparams} collects the full configuration.

\begin{table}[t]
\centering
\caption{DE-MLP architecture and training configuration (production settings).}
\label{tab:hyperparams}
\begin{tabular}{@{}ll@{}}
\toprule
Component & Setting \\
\midrule
Ensemble members $M$        & 5 (seeds $s_0 + m$) \\
Hidden layers               & 64, 32 \\
Hidden block                & Linear $\to$ LayerNorm $\to$ ReLU $\to$ Dropout(0.1) \\
Output head                 & Linear $\to$ 2 (mean, log-variance) \\
Log-variance clamp          & $[-10, 10]$ \\
Loss                        & $\beta$-NLL, $\beta = 0.5$ \\
Optimizer                   & AdamW \\
Learning rate (max)         & $10^{-3}$ \\
Weight decay                & $10^{-4}$ \\
LR schedule                 & one-cycle (\texttt{pct\_start} 0.05, div 25, final div 1000) \\
Gradient clipping           & 1.0 (max norm) \\
Batch size                  & 4096 \\
Maximum epochs              & 500 \\
Early stopping              & patience 50, on the calibration block \\
Calibration fraction        & 0.15 of the training dates \\
Conformal intervals         & Mondrian, 10 predicted-anomaly decile bins \\
\bottomrule
\end{tabular}
\end{table}

% Requires in the main preamble:
%   \usepackage{amsmath}
%   \usepackage{booktabs}
%   \usepackage{natbib}   % \citet  (or map to your journal's citation macro)
%
% NOTE: batch size / max-epoch values below are the production (Linux/HPC) run
% settings; the macOS quick-iteration defaults (2048 / 100) are not used for the
% frozen models.

\section{Deep-ensemble loss, aggregation, and training configuration}
\label{app:model}

\subsection{Member architecture and probabilistic head}

Each of the $M = 5$ ensemble members is a Gaussian-head multilayer perceptron
with two hidden layers (64 and 32 units). Every hidden block applies, in order,
a linear map, layer normalization, a ReLU activation, and 10\% dropout; dropout
is active during training only, so there is no Monte-Carlo dropout at inference
and the ensemble spread alone provides the epistemic uncertainty. For footprint
$i$, member $m$ outputs a predicted mean $\mu_{m,i}$ and a raw second output that
is clamped to $[-10, 10]$ and interpreted as the log-variance
$s_{m,i} = \log\sigma_{m,i}^{2}$, giving a per-footprint predictive variance
$\sigma_{m,i}^{2} = \exp(s_{m,i})$.

\subsection{Training loss}

Writing the target as $y_i \equiv \Delta \mathrm{X_{CO2,\,i}^{B11}}$, the standard
Gaussian negative log-likelihood of one footprint is, up to an additive constant,
%
\begin{equation}
  \ell^{\mathrm{NLL}}_{m,i}
  = \frac{1}{2}\left[\,s_{m,i}
  + \frac{(y_i-\mu_{m,i})^{2}}{\sigma_{m,i}^{2}}\,\right].
  \label{eq:gauss-nll}
\end{equation}
%
Because this loss down-weights high-variance footprints, the mean is poorly
fitted exactly in the near-cloud regime where $\sigma_{m,i}^{2}$ is large. We
therefore train the members with the $\beta$-NLL loss of
\citet{seitzer2022pitfalls}, which multiplies each footprint's NLL by a
stop-gradient function of its own predicted variance,
%
\begin{equation}
  \ell^{\beta\text{-}\mathrm{NLL}}_{m,i}
  = \bigl\lfloor\sigma_{m,i}^{2}\bigr\rfloor_{\mathrm{sg}}^{\,\beta}\;
  \ell^{\mathrm{NLL}}_{m,i},
  \qquad \beta = 0.5,
  \label{eq:beta-nll}
\end{equation}
%
where $\lfloor\cdot\rfloor_{\mathrm{sg}}$ denotes a stop-gradient (the weight is
treated as a constant during backpropagation). Here $\beta = 0$ recovers the
plain NLL of Eq.~\eqref{eq:gauss-nll} and $\beta = 1$ recovers a
mean-squared-error weighting of the mean; the production choice $\beta = 0.5$
restores the mean-fitting gradient in the high-variance near-cloud footprints
while keeping a calibrated variance. The member loss is the average of
Eq.~\eqref{eq:beta-nll} over the training footprints.

\subsection{Ensemble aggregation}

The five members differ only in their random seed (member $m$ uses seed
$s_0 + m$), so their predictions decorrelate and the ensemble spread provides
the epistemic uncertainty. The per-member means and variances are combined as a
Gaussian mixture,
%
\begin{equation}
  \mu^{*}_i = \frac{1}{M}\sum_{m=1}^{M}\mu_{m,i},
  \qquad
  \sigma^{*2}_i
  = \underbrace{\frac{1}{M}\sum_{m=1}^{M}\sigma_{m,i}^{2}}_{\text{aleatoric}}
  + \underbrace{\frac{1}{M}\sum_{m=1}^{M}\mu_{m,i}^{2} - \mu^{*2}_i}_{\text{epistemic}},
  \label{eq:mixture}
\end{equation}
%
which is the law of total variance: the first term is the mean aleatoric
variance predicted by the heads and the second is the between-member (epistemic)
variance. The ensemble mean $\mu^{*}_i$ is the point correction applied to each
footprint.

\subsection{Optimization and leakage control}

Each member is optimized with AdamW (learning rate $10^{-3}$, weight decay
$10^{-4}$) under a one-cycle learning-rate schedule with gradient-norm clipping
at $1.0$, a batch size of $4096$, and up to $500$ epochs. Training is stopped
early on the held-out calibration block (patience of $50$ epochs), and the
checkpoint with the lowest calibration loss is kept; the calibration block is
carved from the training dates and never overlaps the held-out fold or the
validation references. The calibration block is also used to recalibrate the
Mondrian conformal intervals, which are binned by predicted-anomaly deciles.
Feature scaling and any profile-EOF transforms are fitted on the proper-training
partition only. Table~\ref{tab:hyperparams} collects the full configuration.

\begin{table}[t]
\centering
\caption{DE-MLP architecture and training configuration (production settings).}
\label{tab:hyperparams}
\begin{tabular}{@{}ll@{}}
\toprule
Component & Setting \\
\midrule
Ensemble members $M$        & 5 (seeds $s_0 + m$) \\
Hidden layers               & 64, 32 \\
Hidden block                & Linear $\to$ LayerNorm $\to$ ReLU $\to$ Dropout(0.1) \\
Output head                 & Linear $\to$ 2 (mean, log-variance) \\
Log-variance clamp          & $[-10, 10]$ \\
Loss                        & $\beta$-NLL, $\beta = 0.5$ \\
Optimizer                   & AdamW \\
Learning rate (max)         & $10^{-3}$ \\
Weight decay                & $10^{-4}$ \\
LR schedule                 & one-cycle (\texttt{pct\_start} 0.05, div 25, final div 1000) \\
Gradient clipping           & 1.0 (max norm) \\
Batch size                  & 4096 \\
Maximum epochs              & 500 \\
Early stopping              & patience 50, on the calibration block \\
Calibration fraction        & 0.15 of the training dates \\
Conformal intervals         & Mondrian, 10 predicted-anomaly decile bins \\
\bottomrule
\end{tabular}
\end{table}

% Requires in the main preamble:
%   \usepackage{amsmath}
%   \usepackage{booktabs}
%   \usepackage{natbib}   % \citet  (or map to your journal's citation macro)
%
% NOTE: batch size / max-epoch values below are the production (Linux/HPC) run
% settings; the macOS quick-iteration defaults (2048 / 100) are not used for the
% frozen models.

\section{Deep-ensemble loss, aggregation, and training configuration}
\label{app:model}

\subsection{Member architecture and probabilistic head}

Each of the $M = 5$ ensemble members is a Gaussian-head multilayer perceptron
with two hidden layers (64 and 32 units). Every hidden block applies, in order,
a linear map, layer normalization, a ReLU activation, and 10\% dropout; dropout
is active during training only, so there is no Monte-Carlo dropout at inference
and the ensemble spread alone provides the epistemic uncertainty. For footprint
$i$, member $m$ outputs a predicted mean $\mu_{m,i}$ and a raw second output that
is clamped to $[-10, 10]$ and interpreted as the log-variance
$s_{m,i} = \log\sigma_{m,i}^{2}$, giving a per-footprint predictive variance
$\sigma_{m,i}^{2} = \exp(s_{m,i})$.

\subsection{Training loss}

Writing the target as $y_i \equiv \Delta \mathrm{X_{CO2,\,i}^{B11}}$, the standard
Gaussian negative log-likelihood of one footprint is, up to an additive constant,
%
\begin{equation}
  \ell^{\mathrm{NLL}}_{m,i}
  = \frac{1}{2}\left[\,s_{m,i}
  + \frac{(y_i-\mu_{m,i})^{2}}{\sigma_{m,i}^{2}}\,\right].
  \label{eq:gauss-nll}
\end{equation}
%
Because this loss down-weights high-variance footprints, the mean is poorly
fitted exactly in the near-cloud regime where $\sigma_{m,i}^{2}$ is large. We
therefore train the members with the $\beta$-NLL loss of
\citet{seitzer2022pitfalls}, which multiplies each footprint's NLL by a
stop-gradient function of its own predicted variance,
%
\begin{equation}
  \ell^{\beta\text{-}\mathrm{NLL}}_{m,i}
  = \bigl\lfloor\sigma_{m,i}^{2}\bigr\rfloor_{\mathrm{sg}}^{\,\beta}\;
  \ell^{\mathrm{NLL}}_{m,i},
  \qquad \beta = 0.5,
  \label{eq:beta-nll}
\end{equation}
%
where $\lfloor\cdot\rfloor_{\mathrm{sg}}$ denotes a stop-gradient (the weight is
treated as a constant during backpropagation). Here $\beta = 0$ recovers the
plain NLL of Eq.~\eqref{eq:gauss-nll} and $\beta = 1$ recovers a
mean-squared-error weighting of the mean; the production choice $\beta = 0.5$
restores the mean-fitting gradient in the high-variance near-cloud footprints
while keeping a calibrated variance. The member loss is the average of
Eq.~\eqref{eq:beta-nll} over the training footprints.

\subsection{Ensemble aggregation}

The five members differ only in their random seed (member $m$ uses seed
$s_0 + m$), so their predictions decorrelate and the ensemble spread provides
the epistemic uncertainty. The per-member means and variances are combined as a
Gaussian mixture,
%
\begin{equation}
  \mu^{*}_i = \frac{1}{M}\sum_{m=1}^{M}\mu_{m,i},
  \qquad
  \sigma^{*2}_i
  = \underbrace{\frac{1}{M}\sum_{m=1}^{M}\sigma_{m,i}^{2}}_{\text{aleatoric}}
  + \underbrace{\frac{1}{M}\sum_{m=1}^{M}\mu_{m,i}^{2} - \mu^{*2}_i}_{\text{epistemic}},
  \label{eq:mixture}
\end{equation}
%
which is the law of total variance: the first term is the mean aleatoric
variance predicted by the heads and the second is the between-member (epistemic)
variance. The ensemble mean $\mu^{*}_i$ is the point correction applied to each
footprint.

\subsection{Optimization and leakage control}

Each member is optimized with AdamW (learning rate $10^{-3}$, weight decay
$10^{-4}$) under a one-cycle learning-rate schedule with gradient-norm clipping
at $1.0$, a batch size of $4096$, and up to $500$ epochs. Training is stopped
early on the held-out calibration block (patience of $50$ epochs), and the
checkpoint with the lowest calibration loss is kept; the calibration block is
carved from the training dates and never overlaps the held-out fold or the
validation references. The calibration block is also used to recalibrate the
Mondrian conformal intervals, which are binned by predicted-anomaly deciles.
Feature scaling and any profile-EOF transforms are fitted on the proper-training
partition only. Table~\ref{tab:hyperparams} collects the full configuration.

\begin{table}[t]
\centering
\caption{DE-MLP architecture and training configuration (production settings).}
\label{tab:hyperparams}
\begin{tabular}{@{}ll@{}}
\toprule
Component & Setting \\
\midrule
Ensemble members $M$        & 5 (seeds $s_0 + m$) \\
Hidden layers               & 64, 32 \\
Hidden block                & Linear $\to$ LayerNorm $\to$ ReLU $\to$ Dropout(0.1) \\
Output head                 & Linear $\to$ 2 (mean, log-variance) \\
Log-variance clamp          & $[-10, 10]$ \\
Loss                        & $\beta$-NLL, $\beta = 0.5$ \\
Optimizer                   & AdamW \\
Learning rate (max)         & $10^{-3}$ \\
Weight decay                & $10^{-4}$ \\
LR schedule                 & one-cycle (\texttt{pct\_start} 0.05, div 25, final div 1000) \\
Gradient clipping           & 1.0 (max norm) \\
Batch size                  & 4096 \\
Maximum epochs              & 500 \\
Early stopping              & patience 50, on the calibration block \\
Calibration fraction        & 0.15 of the training dates \\
Conformal intervals         & Mondrian, 10 predicted-anomaly decile bins \\
\bottomrule
\end{tabular}
\end{table}
